% !TeX root = ../main.tex

\chapter{Algorithmen}

\section{Motivation}

\dfn{Elementare Addition}{
  Die elementare Addition bezeichnet die Addition dreier einstelliger
  Zahlen, deren Ergebnis höchstens zweistellig ist.
}

\nt{
  Zur Addition zweier $n$-stelliger Zahlen benötigt man im üblichen
  schriftlichen Additionsverfahren $\mathcal{O}(n)$ elementare
  Rechenschritte.

  Für die Multiplikation zweier $n$-stelliger Zahlen ist die Situation
  wesentlich komplexer. Die Laufzeit hängt vom verwendeten Algorithmus
  ab. Moderne Multiplikationsalgorithmen erreichen eine Laufzeit von
  $\mathcal{O}(n\log n)$ bis auf zusätzliche logarithmische Faktoren.
}

\dfn{Elementare Multiplikation}{
  Die elementare Multiplikation bezeichnet die Multiplikation zweier
  einstelliger Zahlen, deren Ergebnis höchstens zweistellig ist.
}

\nt{
  Die Frage, wie schnell sich zwei große Zahlen multiplizieren lassen,
  ist ein zentrales Problem der algorithmischen Mathematik. Es existieren
  Algorithmen, die asymptotisch wesentlich schneller sind als die
  klassische schriftliche Multiplikation.
}

\ex{Numerische Instabilität}{
  Betrachte das lineare Gleichungssystem
  \[
    \begin{aligned}
      10^{-20}x + 2y &= 1, \\
      10^{-20}x + 10^{-20}y &= 10^{-20}.
    \end{aligned}
  \]

  Subtrahiert man die zweite Gleichung von der ersten, erhält man
  \[
    (2-10^{-20})y = 1-10^{-20}.
  \]
  Damit folgt
  \[
    y = \frac{1-10^{-20}}{2-10^{-20}}
      \approx 0{,}5.
  \]

  Wird dieser gerundete Wert in die ursprünglichen Gleichungen
  eingesetzt, können aufgrund der Rundungsfehler widersprüchliche
  Ergebnisse für $x$ entstehen.

  Die exakte Lösung lautet hingegen
  \[
    \begin{aligned}
      y &= \frac{1-10^{-20}}{2-10^{-20}}, \\
      x &= \frac{1}{2-10^{-20}}.
    \end{aligned}
  \]

  Beide Werte liegen sehr nahe bei $0{,}5$. Das Beispiel illustriert,
  dass Rundungsfehler bei numerischen Berechnungen zu erheblichen
  Ungenauigkeiten führen können.
}

\section{Mengen und Relationen}

\dfn{Kartesisches Produkt}{
  Seien $A$ und $B$ Mengen. Das \textbf{kartesische Produkt} von $A$
  und $B$ ist definiert durch
  \[
    A \times B
    := \{(a,b) \mid a \in A,\ b \in B\}.
  \]

  Eine \textbf{Relation} zwischen $A$ und $B$ ist eine Teilmenge
  $R \subseteq A \times B$.

  Für $(a,b) \in R$ schreibt man auch $a\,R\,b$.
}

\ex{Kartesisches Produkt und Relation}{
  Seien
  \[
    A = \{1,3,5\},
    \qquad
    B = \{2,5\}.
  \]
  Dann ist
  \[
    A \times B
    = \{(1,2),(1,5),(3,2),(3,5),(5,2),(5,5)\}.
  \]

  Eine mögliche Relation zwischen $A$ und $B$ ist beispielsweise
  \[
    R = \{(1,2),(3,2),(5,5)\}.
  \]
  Es gilt $R \subseteq A \times B$. Die Relation enthält genau die
  ausgewählten geordneten Paare.
}

\dfn{Abbildung}{
  Seien $A$ und $B$ Mengen. Eine Teilmenge $f \subseteq A \times B$
  heißt \textbf{Funktion} oder \textbf{Abbildung} von $A$ nach $B$,
  falls für jedes $a \in A$ genau ein $b \in B$ existiert, sodass
  $(a,b) \in f$ gilt.

  Schreibweisen:
  \[
    f \colon A \to B, \qquad
    (a,b) \in f \iff f(a)=b, \qquad
    a \mapsto f(a).
  \]
  Die Menge $A$ heißt \textbf{Definitionsbereich}, die Menge $B$
  \textbf{Zielmenge} von $f$.
}

\dfn{Injektivität}{
  Eine Funktion $f \colon A \to B$ heißt \textbf{injektiv}, falls
  für alle $x,y \in A$ gilt:
  \[
    f(x)=f(y) \implies x=y.
  \]
  Äquivalent dazu gilt
  \[
    x \neq y \implies f(x) \neq f(y).
  \]
}

\dfn{Surjektivität}{
  Eine Funktion $f \colon A \to B$ heißt \textbf{surjektiv}, falls
  jedes Element der Zielmenge angenommen wird, das heißt
  \[
    \forall b \in B\ \exists a \in A: f(a)=b.
  \]
}

\dfn{Bijektivität}{
  Eine Funktion $f \colon A \to B$ heißt \textbf{bijektiv}, falls
  sie sowohl injektiv als auch surjektiv ist.
  In diesem Fall wird jedes Element von $B$ von genau einem Element
  aus $A$ angenommen.
}

\ex{Injektivität und Surjektivität}{
  Betrachte die Funktion
  \[
    f \colon \bbN \to \bbN, \qquad f(x)=x^2.
  \]
  Für $a,a' \in \bbN$ gilt
  \[
    a^2=(a')^2 \implies a=a',
  \]
  also ist $f$ injektiv.

  Die Funktion ist jedoch nicht surjektiv, da beispielsweise
  $2 \in \bbN$ kein Urbild besitzt: Es gibt kein $a \in \bbN$
  mit $a^2=2$. Folglich ist $f$ nicht bijektiv.

  Dagegen ist die Funktion
  \[
    f \colon \{0\} \to \{0\}, \qquad f(0)=0,
  \]
  bijektiv.
}

\dfn{Ganzzahlige Intervalle}{
  Für $a,b \in \bbZ$ bezeichne
  \[
    \{a,\dots,b\}
    := \{x \in \bbZ \mid a \leq x \leq b\}.
  \]
  Für $b<a$ setzen wir
  \[
    \{a,\dots,b\} := \emptyset.
  \]
}

\dfn{Endlichkeit}{
  Eine Menge $A$ heißt \textbf{endlich}, falls ein $n \in \bbN_0$
  und eine injektive Funktion
  \[
    f \colon A \to \{1,\dots,n\}
  \]
  existieren.

  Andernfalls heißt $A$ \textbf{unendlich}.
  Für eine endliche Menge $A$ bezeichnet $|A|$ die Anzahl ihrer
  Elemente.
}

\dfn{Abzählbarkeit}{
  Eine Menge $A$ heißt \textbf{abzählbar}, falls eine injektive
  Funktion $f \colon A \to \bbN$ existiert.

  Eine Menge, die nicht abzählbar ist, heißt
  \textbf{überabzählbar}.
}

\dfn{Wörter und Sprachen}{
  Sei $A$ eine endliche, nicht leere Menge, die als
  \textbf{Alphabet} bezeichnet wird. Für $k \in \bbN_0$ sei
  \[
    A^k := \{f \mid f \colon \{1,\dots,k\} \to A\}
  \]
  die Menge aller Wörter der Länge $k$ über dem Alphabet $A$.

  Ein Element $f \in A^k$ wird mit der Folge
  \[
    f(1)f(2)\cdots f(k)
  \]
  identifiziert und als \textbf{Wort der Länge $k$} bezeichnet.

  Die Menge $A^0$ enthält genau ein Element, das
  \textbf{leere Wort} $\varepsilon$.

  Die Menge aller endlichen Wörter über $A$ ist
  \[
    A^* := \bigcup_{k \in \bbN_0} A^k.
  \]
  Eine \textbf{Sprache} über dem Alphabet $A$ ist eine Teilmenge
  von $A^*$.
}

\dfn{Berechnungsproblem}{
  Ein \textbf{Berechnungsproblem} ist eine Relation
  \[
    P \subseteq D \times E,
  \]
  wobei $D$ die Menge der Eingaben und $E$ die Menge der möglichen
  Ausgaben bezeichnet. Es wird vorausgesetzt, dass zu jeder Eingabe
  $d \in D$ mindestens eine Ausgabe $e \in E$ existiert, sodass
  $(d,e) \in P$ gilt.
}

\nt{
  Ein Paar $(d,e) \in P$ beschreibt eine Lösung des Problems für
  die Eingabe $d$. In diesem Fall heißt $e$ eine
  \textbf{korrekte Ausgabe} für die Eingabe $d$.

  Die Elemente von $D$ heißen \textbf{Instanzen} des Problems.
}

\dfn{Eindeutigkeit und Problemtypen}{
  Ein Berechnungsproblem $P \subseteq D \times E$ heißt
  \textbf{eindeutig}, falls zu jeder Eingabe $d \in D$ genau eine
  korrekte Ausgabe $e \in E$ existiert. In diesem Fall beschreibt
  $P$ eine Funktion $P \colon D \to E$.

  Sind $D$ und $E$ Sprachen über einem endlichen Alphabet $A$,
  so heißt $P$ ein \textbf{diskretes Berechnungsproblem}.

  Sind $D \subseteq \bbR^m$ und $E \subseteq \bbR^n$ für
  $m,n \in \bbN$, so heißt $P$ ein \textbf{numerisches
  Berechnungsproblem}.

  Ein eindeutiges Berechnungsproblem $P \colon D \to E$ mit
  $|E|=2$ heißt \textbf{Entscheidungsproblem}.
}

\section{Algorithmen und Programme}

\dfn{Bit}{
  Ein \textbf{Bit} (binary digit) ist eine Informationseinheit,
  die genau einen von zwei Werten annehmen kann: $0$ oder $1$.
}

\ex{Binärdarstellung}{
  Die Zahl $10$ lässt sich im Binärsystem wie folgt darstellen:
  \[
    (1010)_2 = 1\cdot 2^3 + 0\cdot 2^2
             + 1\cdot 2^1 + 0\cdot 2^0 = 10.
  \]
}

\nt{
  Unterschiedliche formale Definitionen von Algorithmen können
  unter geeigneten Voraussetzungen dieselbe Berechenbarkeit
  beschreiben.
}

\nt{
  Ein typischer Kompilierungsprozess lässt sich vereinfacht
  folgendermaßen darstellen:
  \[
    \text{Quellcode}
    \longrightarrow
    \text{Maschinencode}.
  \]
  Dazwischen können weitere Darstellungen liegen, beispielsweise
  Zwischencode oder eine interne Repräsentation.
}

\dfn{Programm}{
  Ein \textbf{Programm} ist eine Zeichenfolge, die den syntaktischen
  Regeln einer Programmiersprache genügt. Seine Ausführung wird
  durch die Semantik der jeweiligen Programmiersprache bestimmt.
}
